\documentclass[11pt,a4paper]{ctexart}
\usepackage[margin=2.4cm]{geometry}
\usepackage{amsmath,amssymb}
\usepackage{fancyhdr}
\usepackage{enumitem}
\linespread{1.18}
\setlength{\headheight}{15pt}
\setlength{\parskip}{0pt}

\pagestyle{fancy}
\fancyhf{}
\fancyhead[L]{\small Team 2322687}
\fancyhead[R]{\small \thepage}
\renewcommand{\headrulewidth}{0.4pt}
\fancypagestyle{plain}{\fancyhf{}\renewcommand{\headrulewidth}{0pt}}

\ctexset{
  section/format = \Large\bfseries\raggedright,
  subsection/format = \large\bfseries\raggedright,
  subsubsection/format = \normalsize\bfseries\raggedright
}
\setcounter{secnumdepth}{3}

\begin{document}

%====================== 摘要页 ======================
\thispagestyle{empty}
\begin{center}
{\small 2023 MCM/ICM\quad 选题 A\quad 队伍控制号 2322687\quad 摘要页（Summary Sheet）}\\[18pt]
{\LARGE\bfseries 还有空间吗？植物群落的时空偏微分模型}\\[8pt]
{\large （Is There Space for More? A Spatiotemporal Partial Differential Equations Model for Plant Growth During Drought Cycles）}\\[10pt]
\end{center}

\begin{center}
{\large\bfseries 摘\quad 要}
\end{center}

\vspace{6pt}
干旱对植物群落的存续构成重大风险。尽管单个植物个体可能难以适应干旱所造成的水分匮乏，但既有文献提示，植物群落的生物多样性与它在干旱周期中存活的能力之间存在正相关关系。我们受命验证这一论断，并建立一个与生物群落类型无关的、具有普适性的模型，用以预测植物群落在干旱周期中的变化。

我们的做法是建立一个参数化的偏微分方程组模型，依据可用水量刻画多个植物物种随时间生长的过程。关键在于，我们还将空间维度与空间分析纳入模型，从而能够考虑植物扩散以及水分跨越物理边界的扩散。为支撑这一生长模型，我们还实现了一个降雨子模型，该子模型能够通过改变干旱发生的时间与严重程度来模拟多种不同类型的干旱。将该子模型并入植物生长模型之后，我们使用 Mathematica 对偏微分方程进行数值分析，从而模拟植物生长。我们对模型参数进行了灵敏度分析，通过从一个物种的对照组出发、逐步增加独特物种的数目来模拟生物多样性的效应。我们还分析了不同类型干旱的影响，以及按抗性与恢复力衡量的不同植物类型的影响。最后，我们研究了污染或栖息地缩减等外部因素的影响。

最终我们发现：尽管在给定群落中增加独特植物的数目会引发种间竞争，但它通过提高长期植物生物量的总量，累积性地使整个系统受益。然而，边际收益随独特物种总数的增加而递减：从一个物种增加到两个物种会在总生物量上产生显著收益，但从两个增加到三个时变化较小，而从三个增加到四个实际上会降低总生物量。与此同时，我们的灵敏度分析发现，恢复力与抗性的变化会相应地影响总植物生物量，但并不影响模型的总体行为。此外，我们确定栖息地破坏会因空间竞争的加剧而使个别物种在植物群落中占据主导地位。我们发现，在某些极端条件下，特定植物物种究竟会长期存活并占据主导，还是会走向灭绝，取决于其初始条件。

我们这一独特的空间方法，相对于只依赖时间的传统、与空间无关的模型（如经典的 Lotka-Volterra 模型）而言，其主要优势在于它能够刻画植物与地下水的基本物理行为——即在二维空间中的扩散与弥散。若不考虑空间因素，便无法准确刻画这样一个事实：一般而言，植物与水分只能传播到相邻位置。这也使我们能够对诸如在给定位置上争夺栖息空间或水分的现象进行建模。最后，空间方法还提供了模拟随机初始位置（从而镜像真实世界环境的随机性）以及便捷地预测栖息地缩减影响的能力。结合一个基于真实降雨数据与干旱条件的模型，我们的模型可作为模拟任何给定群落在干旱周期下植物生长的综合性、普适性的唯象模型。

\vspace{6pt}
\noindent\textbf{关键词：}偏微分方程，时空分析，干旱建模

\newpage
\pagestyle{fancy}

%%====================== 正文 ======================
\section{引言}

\subsection{背景}

研究干旱周期如何影响植物群落，对于理解气候变化对生态系统与农业的影响至关重要。干旱能够显著改变植物的生长、分布与生产力，进而对生态系统功能与服务（如养分循环、土壤稳定性与碳储存）产生连锁效应。此外，干旱还会影响农业生产力，因为它们降低作物产量、增加灌溉需求，并加剧土壤侵蚀与土地退化 [1]。通过研究干旱如何影响植物群落，科学家能够制定有效的管理策略以缓解干旱对生态系统的冲击，这包括培育抗旱作物品种、提高水分利用效率，以及实施可持续的土地管理实践。归根结底，在气候变化背景下，理解干旱对植物群落的影响对于维持生态系统健康、支撑可持续农业至关重要。

植物物种常常在地理上被归并为群落，以便更好地理解物种之间的相互作用与关系 [2]。群落内各物种的相对组成并非静止不变，而是随着对外部刺激与胁迫的响应而改变。特别是，能够剧烈改变群落组成的一种常见胁迫来源是干旱期 [3]。

由于不同生物具有不同特性，它们对极端条件的反应与适应方式也各不相同。此类响应可由恢复力（resilience，即在极端条件之后恢复的能力）与抗性（resistance，即在极端条件下仍能存续并生长的能力）这两个概念很好地刻画 [4]。既有研究表明，相对于更为同质的群落，这些群落内部的局部生物多样性会增强群落整体适应干旱期的能力。有研究提出一种可能性：多样化群落所提供的遗传多样性，使物种多样的群落在长时间少雨甚至无雨的时期也更有可能有效繁衍 [5]。然而，关于生物多样性在多大程度上（定量地）增强群落在干旱中存续与生长的能力，以及这种存续能力如何随物种的具体数目而变化，目前研究甚少。

\subsection{问题重述}

为回答这一问题，我们团队受命建立一个模型，以预测植物群落的组成如何随干旱周期而变化。该模型及其分析必须：

\begin{enumerate}
\item 考虑物种之间的相互作用，既包括竞争，也包括相互支持，且须在干旱期间成立；
\item 评估需要多少物种才能从生物多样性现象中获益，并评估群落内物种数目增加所带来的影响；
\item 考虑干旱频率与严重程度变化的影响；
\item 考虑群落内不同物种类型的影响；
\item 考虑污染与栖息地破坏的影响。
\end{enumerate}

\subsection{已有文献}

在建立我们自己的定制模型之前，我们研究并分析了若干已有模型。下面给出两个重要的例子。

首先，我们考察了经典的 Lotka-Volterra 模型，作为考虑植物物种之间竞争及其对生长影响的出发点 [6]。捕食者—被捕食者方程组的一个抽象化版本为：
\[
\frac{dx}{dt} = \alpha x - \beta xy , \qquad \frac{dy}{dt} = \delta xy - \phi y
\]
其中 $\alpha$ 与 $\delta$ 分别为被捕食者物种与捕食者物种的增长率，而 $\beta$ 与 $\delta$ 表示两物种之间的竞争性相互作用。

尽管 Lotka-Volterra 模型提供了纳入竞争因素的一种方法，但我们认为它针对捕食者—被捕食者关系的特殊性，使其难以准确刻画植物之间的相互作用，因为这些相互作用也可能具有互利性。与此同时，对本任务而言该模型或许最重要的问题在于其增长率恒定的假设。它没有考虑增长率与水资源稀缺程度之间的关系。

为处理植物建模中资源利用这一方面，我们还考察了用于刻画微生物生长速率的 Monod 方程 [7]：
\[
\mu = \mu_{max}\frac{S}{K_S + S}
\]
其中增长率 $\mu$ 由最大增长率与一个分数项的乘积决定，$S$ 表示限制性底物（即资源）的浓度，$K$ 为常数。

这两个模型为我们的工作提供了基础性的起点，但最终不足以处理手头的问题。特别是，我们希望纳入以下在这些体系中没有被考虑的因素：

\begin{enumerate}
\item 多物种相互作用；
\item 地下水的储存；
\item 群落的空间维度，它影响植物与水分扩散；
\item 不同类型的干旱；
\item 植物的恢复力与抗性；
\item 污染与栖息地缩减等外部因素。
\end{enumerate}

\section{建立植物群落模型}

\subsection{假设}

为简化建模，我们确立若干基本假设，并给出其合理性说明：

\begin{itemize}
\item 假设：植物生物量仅通过植物生长而增加。
  \begin{itemize}
  \item 理由：题目陈述围绕分析植物相对于天气条件的生长而展开。因此，我们认为只需讨论植物生物量如何随植物生长或死亡而变化，而忽略人类开发、捕食者等其他外部因素。
  \end{itemize}
\item 假设：在给定的某次模拟过程中，环境中不会突然出现新的植物物种。
  \begin{itemize}
  \item 理由：与上述假设类似，我们不考察可能使某一植物在生态系统中突然出现的外部影响。相反，我们将把分析集中于生态系统中已有的植物，研究它们的生物量如何随时间变化。
  \end{itemize}
\item 假设：植物仅通过其抗性（多容易死亡）与恢复力（多容易生长）加以区分。
  \begin{itemize}
  \item 理由：在分析极端天气条件下植物生长的语境中，需要考虑的最重要因素是植物在这些条件下如何存活。因此，我们不考察植物结构、形状或大小等因素，而假定给定植物的抗性与恢复力足以概括其其他决定性特征。
  \end{itemize}
\end{itemize}

\subsection{模型概览与论证}

为探讨这一问题，我们的模型将尝试预测地下水量以及每个植物物种的总生物量将如何变化——这些量将作为因变量。具体而言，由于我们关注的是研究这些量的变化率，我们将建立一个基于微分方程的模型，因为微分方程能够自然地刻画随时间的变化。

然而我们主张，仅跨越时间对系统进行分析是不够的，因为生态群落的物理维度对其发展具有高度影响。其一，与针对动物或微生物的模型不同，植物一般不会移动，这因而限制了它们的水分消耗以及在非相邻位置上与其他植物的相互作用。例如，不同植物物种可能由于某一特定位置的水分短缺而无法占据栖息地中的同一物理位置，并且栖息地的大小也会限制植物生长。

此外，水分的散布在很大程度上由其跨空间的扩散以及浓度梯度所驱动。对给定区域内水分的空间运动进行建模，也使我们能够采用对给定位置上水量的更准确度量，而不是系统中的总水量——后者对植物生长的刻画精度要低得多。例如，给定位置上植物的生长并不取决于整个系统的水量，而取决于它所占据的那个特定位置的水量。

因此，我们的模型除时间之外还将把空间维度视为自变量，这指向了偏微分方程而非常微分方程的必要性；这也使我们能够在后续的灵敏度分析中研究栖息地缩减的影响。

在模型内部，我们将纳入若干参数，以便探讨并刻画以下条件：

\begin{itemize}
\item 独特植物物种的数目；
\item 每个物种的恢复力与抗性；
\item 每周的降雨量；
\item 干旱的频率与量级。
\end{itemize}

需要注意，总体而言，除自变量时间以外，模型是无量纲的；也就是说，我们并不度量生物量、地下水、降雨等的具体数量。这样做是因为本研究的目的是确定不同因素对植物生长的更宽泛的定性影响与相对影响，而不是在某一组精确参数值下预测某一特定区域内存在的植物的精确数量。后一种考察恐怕无法用一个唯象的、机理性的数学模型来完成，而恰恰是这样一个模型才能对植物与其群落相互作用的性质与行为给出深刻洞察，这也是本研究的最终目标以及整个题目陈述的目标。因此，我们认为，建立一个可适配于任何给定植物群落的无量纲模型要有意义得多，其结果能够提供关于物种在不同条件下（无论是恢复力与抗性这样的内部条件，还是干旱、污染或栖息地缩减这样的外部条件）如何相互作用的基础性信息。自变量时间被赋予量纲，是为了进行参数值计算，并为我们的发现提供一些基本语境；但即便是它也可以被无量纲化，而不改变模型的底层有效性。

\subsection{变量与记号}

在模型中，我们将引入若干变量与参数。此处给出一个汇总：

\begin{itemize}
\item 物种数目，记号为 $k$：系统中独特物种的数目；
\item 物种生物量，记号为 $n_i$：物种 $i$ 的植物生物量总量；
\item 总植物生物量，记号为 $N$：所有物种的植物生物量总量；
\item 水分，记号为 $w$：地下水总量；
\item 空间维度，记号为 $x,\ y$：所建模的二维空间；
\item 时间，记号为 $t$：时间，以周为单位度量；
\item 植物生长率，记号为 $c_i$：物种 $i$ 的生长率；
\item 植物死亡率，记号为 $m_i$：物种 $i$ 的死亡率；
\item 由浓度差驱动的水流，记号为 $v$：水分由高浓度向低浓度运动的速度；
\item 由扩散驱动的水流，记号为 $d$：水分在平坦地形上流动的速度；
\item 降雨，记号为 $a_t$：时刻 $t$ 的每周降雨量；
\item 栖息地大小，记号为 $L$：（方形）栖息地的长度与宽度。
\end{itemize}

\subsection{植物生物量}

我们首先推导一个刻画每个物种 $i$ 总生物量的方程，以 $n_i(x,y,t)$ 表示，它可受三个因素影响：植物生长、植物死亡，以及植物在栖息地内的扩散。

植物生长由给定物种的增长率、可用水量以及当前植物数量决定。文献表明植物生长并非线性，且给定区域内同种植物的大量存在会促进进一步的植物生长 [8]。因此，我们提出一个相互作用项，它由植物生长率、当前植物生物量的二次型以及水分可用性共同定义。植物生物量仅通过植物死亡而损失，死亡取决于当前植物生物量与物种的死亡率 $m_i$。最后，植物扩散可由跨越 $x$ 与 $y$ 的一般扩散方程刻画。综合起来，这给出如下偏微分方程：
\begin{equation}
\frac{\partial n_i}{\partial t} = c_i w n_i^2 - m n_i + \frac{\partial^2 n_i}{\partial x^2} + \frac{\partial^2 n_i}{\partial y^2}
\end{equation}

\subsection{水与降雨}

我们接着推导一个刻画地下水量 $w$ 的方程，它受降雨、排水（其中包含蒸发）以及水分在栖息地内的散布影响。水分通过随时间变化的降雨 $a_t$ 进入系统，并以与当前地下水量成比例的速率因排水而损失。水分还会因植物消耗而损失，这一项同样应由总植物生物量与水分可用性的二次相互作用来刻画，其依据是生物量越高越能促进未来生长的观念。

关于散布，水分以速度 $v$ 由高浓度流向低浓度，从而引入 $v$ 与 $w$ 的梯度之间的一个相互作用项。水分还可以以速度 $d$ 在平地上流动，从而引入 $w$ 的一个带尺度因子的扩散方程。因此，我们提出：
\begin{equation}
\frac{\partial w}{\partial t} = a_t - w - wN^2 + v\left(\frac{\partial w}{\partial x} + \frac{\partial w}{\partial y}\right) + d\left(\frac{\partial^2 w}{\partial x^2} + \frac{\partial^2 w}{\partial y^2}\right)
\end{equation}

\subsection{物种竞争}

我们模型的一个关键组成部分是研究系统中多个独特物种之间相互作用的能力。正因如此，我们把总植物生物量 $N$ 拆分为每个植物物种 $i$ 的生物量，这使我们能够为每个植物物种写出单独的方程，并把总生物量写为这些方程之和。在数学上，这给出如下系统：
\begin{equation}
\begin{array}{rcl}
\frac{\partial n_1}{\partial t} & = & c_1 w n_1^2 - m n_1 + \frac{\partial^2 n_1}{\partial x^2} + \frac{\partial^2 n_1}{\partial y^2}\\[1.0em]
& \vdots & \\[0.4em]
\frac{\partial n_k}{\partial t} & = & c_k w n_k^2 - m n_k + \frac{\partial^2 n_k}{\partial x^2} + \frac{\partial^2 n_k}{\partial y^2}\\[1.0em]
\frac{\partial N}{\partial t} & = & \frac{\partial n_1}{\partial t} + \cdots + \frac{\partial n_k}{\partial t}\\[1.0em]
\frac{\partial w}{\partial t} & = & a_t - w - wN^2 + v\left(\frac{\partial w}{\partial x} + \frac{\partial w}{\partial y}\right) + d\left(\frac{\partial^2 w}{\partial x^2} + \frac{\partial^2 w}{\partial y^2}\right)
\end{array}
\end{equation}

为研究独特物种如何相互作用、竞争、协作，并最终影响植物群落的干旱适应能力这一现象，我们可以增加物种总数 $k$。

\subsection{降雨、干旱、物种类型与污染}

通过随时间变化的参数 $a_t$，我们可以改变系统在给定时间段内所获得的降雨量。借此，我们可以考虑干旱期（即降雨减少），并调整此类干旱的长度、频率与量级。最后，通过分别控制植物恢复力与抗性的参数 $c_i$ 与 $m_i$，我们能够模拟具有不同属性的不同植物物种类型的影响，以及不同污染程度的影响——污染会影响植物的生长能力与死亡倾向。我们将在模型分析过程中对这些参数进行灵敏度分析。

\subsection{初始条件与边界条件}

与所有偏微分方程组一样，我们还必须确立初始条件与边界条件。每个植物物种的初始生物量取为栖息地内某一随机位置处的值 $1$，代表模拟开始时随机散布的一小簇种子或孢子。地下水的初始水量取为 $t=0$ 时刻的降雨量。

我们将对每个植物物种的生物量与地下水均施加周期边界条件。这在数学建模中是一种合理且常用的条件设定，因为我们试图用一个单元胞来近似一个大型的半无限系统。在拓扑上，由二维周期边界条件所构成的空间可以看作被映射到一个紧空间上。综合起来，并回顾参数 $L$，这给出：
\begin{equation}
\begin{array}{rcl}
n_1(x,y,0) & = & n_1^0(x,y)\\
& \vdots & \\
n_k(x,y,0) & = & n_k^0(x,y)\\
N(x,y,0) & = & n_1^0(x,y) + \cdots + n_k^0(x,y)\\
w(x,y,0) & = & a_0
\end{array}
\end{equation}
作为初始条件，以及
\begin{equation}
\begin{array}{rcl}
n_1(0,y,t) & = & n_1(L,y,t), \quad n_1(x,0,t) = n_1(x,L,t)\\
& \vdots & \\
n_k(0,y,t) & = & n_k(L,y,t), \quad n_k(x,0,t) = n_k(x,L,t)\\
N(0,y,t) & = & N(L,y,t), \quad N(x,0,t) = N(x,L,t)\\
w(0,y,t) & = & w(L,y,t), \quad w(x,0,t) = w(x,L,t)
\end{array}
\end{equation}
作为边界条件。

\subsection{模型分析}

由于其高度非线性的性质，所推导的偏微分方程没有解析解。因此，为研究式（3）所给出的系统，我们在 Mathematica 中对偏微分方程进行数值分析。具体而言，我们使用数值直线法（method of lines），它通过在除一个维度以外的所有维度上对格式进行离散化来求解偏微分方程，随后把半离散问题作为常微分方程组（ODEs）或微分代数方程组（DAEs）进行积分。

\section{降雨与干旱模型}

\subsection{一般降雨模型}

在植物群落模型之外，我们还建立了一个降雨模型，用以模拟自然降雨以及干旱期与恢复期。为此，我们从一个具有平均降雨概率与平均降雨量的一般降雨模型出发。先对降雨概率作均匀随机化，以确定逐周的降雨可能性，随后再对平均降雨量施加另一层随机化。

接着，我们用季节变异指数（它代表季节变动的幅度）来修正这两个数值。例如，某些群落可能经历季节性季风，而另一些群落（如沙漠中的群落）的季节变异则可能极小。由于季节具有周期性，我们将这一季节变异施加于正弦函数之上，以恰当地构造作用于每周降雨量的季节系数：
\[
V^{*}\left(\sin\left(\frac{2\pi}{52}\left(d \bmod 52\right)\right) + \frac{51}{52}\right)
\]
其中 $V$ 为季节变异指数，$d$ 为周数，且 $\sin\left(\frac{2\pi}{52}x\right)$ 的周期为 52 周。

\subsection{干旱调整模型}

在建立降雨模型之后，我们随即在一个 1000 周的一般降雨模型之上叠加干旱期。有些干旱还被刻意安排在雨季之内，以回应题目陈述中“在降水充沛时期也包含干旱”的要求。在 1000 周的数据中总共包含三段相互独立的干旱：具体而言，干旱分别发生在 $t=200$ 至 $t=250$、$t=500$ 至 $t=600$，以及 $t=750$ 至 $t=900$。

我们还设置了参数，用以限制干旱持续时间的范围以及严重程度（即降雨发生时降雨量的减少百分比）。凭借干旱出现时机、持续时间与严重程度这三个参数，我们因而能够生成各自独特的“干旱情景”，这比仅使用一律漫长而严酷的干旱提供了更为全面的干旱刻画。

在确定参数之后，我们随即把设定好的干旱模型施加于一个经随机化的千周降雨时段之上，并与一般降雨模型配合使用。我们在第 4.2 节中分析由不同参数衍生出的不同干旱情景的影响。

降雨与干旱的构造过程可用一个逐周的决策树加以表示：第一层圆圈代表改变下一层条件（是否会有降雨）的修正量，每一次迭代代表一周的决策树，该循环共执行 1000 周。

\section{结果}

\subsection{独特植物物种的数目}

本研究的一个主要目的是考察不同物种在干旱周期中的相互作用，而二者之间无疑存在关联 [9, 10, 11]。因此，我们首先研究系统如何随独特物种数目的不同而变化，从 $k=1$（即仅含一种植物的植物群落）开始，它既可作为对照组，也可用于验证模型的性能。

首先考察单物种情形在若干时刻的空间格局。在 $t=200$（第一次干旱之前）时，植物在栖息地内似乎正在形成条带状的植被；在 $t=750$（第三次干旱之前）时，植物仍然形成条带，但每一条带都更厚、更集中；在 $t=900$（第三次干旱之后）时，干旱过后植物簇变得更小、集中度更低；在 $t=1000$（自干旱中恢复）时，植物能够在干旱之后实现完全恢复。

更系统地说，这些空间密度图给出了植物群落在某一给定时点的空间密度：对栖息地内任一给定的空间坐标，该位置的颜色代表该指定时刻该给定位置上植物的生物量。与之相配的总植物生物量曲线则给出 1000 周期间总植物生物量的变化。可以看到，生物量突然下降的时段与图 2 所示的干旱时段相对应。

这些结果验证了模型的有效性。事实上，所呈现的空间格局与仅含单一物种的群落中真实世界的植物形态相吻合：例如，南非的规则斑块状植被形态与 $t=200$ 及 $t=900$ 时的形态相互映照，而尼日尔的条带状植被形态则与 $t=750$ 及 $t=1000$ 时的形态相互映照 [12]。

在这一运作良好的模型基础上，我们随后把模型扩展到包含两个独特物种（$k=2$）。所有快照均在 $t=1000$ 时取得。

对两个物种的行为给出了引人入胜的洞察：它们似乎直接争夺空间。事实上，在任一给定位置上，要么第一个物种繁盛，要么第二个物种繁盛，但绝不会两者同时繁盛——这在模型的定义下是合理的，因为长期来看，在任一给定的空间坐标上最终只有一个植物物种能够利用该处的水分。

把两物种情形的总生物量曲线与单物种情形相对照可以看出，尽管存在物种竞争，总生物量仍显著增加——实际上增加到两倍以上。这表明物种协作的存在，其机理可以解释为：系统中植物越多，越能通过优化水分利用而保存水分；也就是说，由于水分被保存在植物体内，允许通过排水离开系统的水分就更少。此外还可以观察到，在两物种情景中，植物在物理上占据了整个系统，而单物种情景中的植物则形成簇与条带，未能占据整个栖息地。这与多样化的过程相呼应：植物物种越多样，栖息地内能被占据的空间领域就越多。

接下来我们将探讨当独特植物物种数目继续上升时，这些结果是否依然成立。特别是，我们希望研究：当在任一给定时刻、任一空间坐标上都同时存在多种针对空间与水分的竞争性相互作用时，物种竞争的格局如何改变与演化。因此，我们首先把系统扩展到 $k=3$，即模拟含三个独特物种的植物群落。所有快照均在 $t=1000$ 时取得。

与 $k=2$ 的情形（在任一给定空间坐标上只有一种植物能够繁盛）不同，在 $k=3$ 的情形中物种竞争不那么容易直接观察到。然而，三个物种彼此争夺水分的现象仍可辨识：例如，考察大致由 $((0,100),(100,200))$ 所界定的小区域，占据栖息地大部分区域的物种 3 在该处无法生长；在该区域内部，上部被物种 1 占据，而左下部分被物种 2 占据——只是边界不像两物种情景中那样清晰。此外还可以观察到，在任一给定的空间坐标上，并非三个物种都能繁盛。

然而，我们再次关注以总植物生物量衡量的整个系统的成败与增长。可以看到，尽管存在种间竞争，系统中的生物量总量再次增加——不过其增幅明显小于从一个物种增加到两个物种时的增幅。事实上，此时总生物量的峰值比两植物情景在干旱之间的峰值高出约 20\%，这一增幅虽显著，却远不及从一种植物增加到两种植物时那样显著。这表明系统并不会从更多生物多样性中获得无止境的收益，并且随着独特物种数目的上升，由此产生的种间竞争加剧会对系统的总生物量造成越来越不利的影响。

为进一步考察这一现象，我们用 $k=4$ 进行实验。需要注意，此时的模拟大约从三分之一处、即 $t=350$ 开始。这样做是为了节省计算时间，因为对由五个偏微分方程组成的系统进行数值求解可能极其耗时。所有快照均在 $t=1000$ 时取得。

这一情景还就种间竞争的性质提出了引人入胜的洞察。在四物种情形中，模拟结束时的总生物量低于三物种情形，这表明独特植物物种的数目不可能无止境地增加而使整个系统受益。事实上，在某一临界点上，植物物种之间为争夺水分而产生的竞争量会抵消掉因植物存在而在系统内被保存下来的水分总量。这一阈值未必就是四个物种，它可能取决于随机初始条件与参数；然而这一洞察揭示了一个事实：从唯象层面看，无限增加独特植物物种的数目未必有益于整个系统。

\subsection{不同类型的干旱}

我们还可以考察影响系统的不同类型的干旱；具体而言，我们考察其严重程度（降雨下降有多显著）、持续时间（干旱有多长）以及频率（模拟期间发生多少次干旱）。为确保是干旱的这些性质在影响模型结果，我们将保持其他所有参数不变。为在仍然考虑种间竞争效应的同时减少计算时间，我们选取 $k=2$。

首先，我们考察一种比最初所用干旱更为严酷的干旱，其中干旱期内的降雨水平甚至低于图 2 所示的情形。这里给出两个物种在若干时刻的空间状态：物种 1 在 $t=200$（第一次干旱之前）、$t=600$（第二次干旱之后）、$t=750$（第三次干旱之前）以及 $t=1000$（自干旱中恢复）时的分布；物种 2 在 $t=200$（第一次干旱之前）、$t=750$（第三次干旱之前）、$t=600$（第二次干旱之后）以及 $t=1000$（自干旱中恢复）时的分布；此外还有系统总生物量随时间的变化曲线。

干旱严重程度的提高所带来的影响是：在干旱期之后几乎不再有植物残留，尽管植物成簇的形状仍然保留。然而，植物能够从这些干旱期中迅速恢复，因为它们并未完全灭绝，这一点在物种 2 的各个快照中表现得尤为清楚。同时，也正是严重干旱导致了剧烈的峰谷变化，系统总生物量在某些时点几乎降至零。

接下来，我们考察频率更高的干旱，其中干旱的严重程度较低，但发生得更为频繁。具体而言，我们在图 2 所示降雨模型的基础上加以扩展，纳入 $t=200$ 至 $t=300$、$t=400$ 至 $t=500$、$t=600$ 至 $t=700$，以及 $t=800$ 至 $t=900$ 的若干段干旱。此处同样给出两个物种在若干时刻的空间状态：物种 1 在 $t=200$（第一次干旱之前）、$t=650$（第二次干旱之后）以及 $t=1000$（自干旱中恢复）时的分布；物种 2 在 $t=200$（第一次干旱之前）、$t=650$（第二次干旱之后）以及 $t=1000$（自干旱中恢复）时的分布；此外还有系统总生物量随时间的变化曲线。

与更为严酷的干旱不同，更频繁的干旱从不曾威胁到植物群落的整体存在。事实上，在干旱期之后植物依然相当活跃，而干旱期与非干旱期之间的差异也不那么极端。这表明，一般而言，更频繁但较不严酷的干旱要优于其反面情形。

\section{灵敏度分析}

我们还有兴趣研究系统中植物物种的类型，以及污染与栖息地缩减等其他因素。

\subsection{物种类型}

生长率与死亡率的参数化使我们能够研究不同类型的植物——具体而言，我们聚焦于生态学分类中常见的两个主要特征：恢复力与抗性。恢复力被定义为某一植物作为一个物种在干旱之后恢复的能力，即其通过繁殖恢复到先前生物量的能力。抗性被定义为某一物种中的单个植物在环境胁迫下存活的能力。

因此，我们可以用生长率参数 $c_i$ 作为衡量植物物种恢复力影响的代理量。具体而言，回顾基线值 $c_i=1$，我们考察两种情形：$c_i=0.8$（植物恢复力较低）与 $c_i=1.2$（植物恢复力较高）。所有快照均在 $t=1000$ 时取得。

引入恢复力较低（即生长率较低）的植物物种，所产生的结果与最初的模拟相似。然而正如预期的那样，可以看到每个植物物种的总体生物量都显著降低。特别地，尽管生长因子仅下降了 20\%，系统总体生物量的变化幅度却显著更大，这表明植物恢复力对增长的驱动并非仅呈线性。对照模型本身来看这是合理的：该偏微分方程组由若干相互作用项构成，这些项在参数与因变量之间生成了复杂的非线性关系。

与此同时，引入恢复力更高（生长率更高）的植物会产生相应的影响。虽然系统的总体行为与原始示例相对一致，但每个物种的总体生物量都显著更高。同样，由此产生的生物量增幅大于恢复力本身的增幅，表明模型的强互联性产生了非线性关系。具体而言，随着植物生物量增加，系统中被保存的水分更多，这意味着因蒸发而损失到系统之外的水分更少。

\subsection{抗性与污染}

继续分析植物物种抗性变化的影响，我们可以用植物死亡率 $m_i$ 作为抗性的代理量。从逻辑上讲，在相似的环境条件下，抗性较高的植物相比抗性较低的植物具有更低的死亡率。

我们的模型还能够刻画系统中污染的存在所引起的变化。我们知道，在相同水分消耗水平下，受污染的植物表现出更高的死亡率 [13]。事实上，衡量植物应对胁迫能力的恢复力特征同样可以适用于污染这一环境胁迫。

因此，我们通过改变植物物种的死亡率 $m_i$，在植物抗性的基础上模拟全栖息地尺度污染的有无所造成的影响。具体而言，给定基线值 $m=0.05$，我们模拟了 $m=0.03$（污染较轻）与 $m=0.07$（污染较重）两种情景。所有快照均在 $t=1000$ 时取得。

考察减轻污染、从而降低死亡率所带来的后果：正如预期，总植物生物量随时间增加，而系统的总体行为仍保持相似。植物的成簇格局与此前的发现相对一致。

类似地，考察提高污染水平、从而推高植物死亡率的后果。在这种情况下，植物生物量在所有时段都下降，但系统的总体行为再次保持一致。请注意，在两种情景中，系统总生物量的增减幅度都大于死亡率本身的变化幅度，这强调了偏微分方程系统的互联性如何等价于对参数值的非线性响应。

\subsection{栖息地缩减}

栖息地缩减是建模植物生长时应当考虑的另一个关键因素，在当代气候条件下它变得日益重要。得益于模型的空间能力，我们可以通过减小可供植物利用的总区域（由 $L$ 支配）来便捷地刻画栖息地缩减。

我们考察两种可能的缩减：一种是幅度较小的缩减，取 $L=300$；另一种是幅度较大的缩减，取 $L=100$。

当栖息地经历幅度较小的规模缩减时，植物生长的行为看起来相似，只是独立簇的数目更少。这是合理的，因为可供植物生长的总面积减少自然会减少它们能够形成的簇的数目。然而，系统的整体动态与原始情形看起来相似；唯一的差别是总生物量如预期那样下降了。

幅度较大的栖息地缩减给出了另一组极为有趣的结果。事实上，当给定栖息地小到一定程度时，一种植物会完全死亡，而另一种则得以存活并扩散至整个系统。这进一步支持了物种之间空间竞争的存在，而当栖息地剧烈缩减时这种竞争被凸显出来。可以看到，第一个物种随时间逐渐死亡，即便在降雨充足的时段也无法恢复生机；而第二个物种则完全接管了整个系统，这大概是由于其初始条件更优。

这一双物种的试验揭示了栖息地破坏的一个至关重要的总体负面效应：存在这样一个临界点，更小的空间栖息地由于对有限资源的竞争加剧而无法再支撑两个物种，从而导致两个物种之一的灭绝。这一将栖息地缩减与生物多样性下降联系起来的观念可以推广到任意数目的物种。因此，把栖息地破坏推向极端，其最终结果可以说是只留下一个存活物种，它在与其他物种的竞争中最为出色，但重要的是，它未必最擅长在干旱等环境胁迫下存活。

因此，当把栖息地破坏与干旱周期一并考虑时，我们从植物群落模型中看到，由栖息地破坏所导致物种数目的减少意味着该群落对干旱变得更加脆弱、更不稳定。

\section{优点与不足}

\subsection{优点}

\begin{itemize}
\item 空间模型的使用具有动态性与全面性。我们不仅考虑降雨，还考虑地下水的扩散，从而刻画了地表的自然渗透性。位置相关的竞争契合了不可移动植物在空间上局域化的事实，以及真实世界干旱情景中植物自然成簇的现象。
\item 两个模型中内建的随机化模拟了真实世界。随机的初始播种带来更自然的初始植物散布。把每日降雨在特定范围内随机化，比使用恒定的降雨函数更贴近自然。
\item 我们参数化的降雨与干旱模型能够模拟种类广泛的不同干旱。在长度、严重程度与出现时机上的可选性，使得在分析干旱周期对植物群落模型的影响时具有更好的客观性与可控性。这与真实世界数据形成对照：真实数据全部是观测性的，可能不完整，且不像我们这样受控，从而难以有效地分离并检验这三个参数。
\item 建立唯象的、机理性的模型（例如基于变量之间理论关系的模型）而非数据驱动的模型，使我们能够刻画系统的互联性，并稳健地评估植物群落各个方面的关系。
\end{itemize}

\subsection{不足}

\begin{itemize}
\item 我们的模型需要大量计算能力，因为对 1000 周中的每一周，模型都要在 $500\times 500$ 的空间网格上逐点考虑水分可用性与植物竞争。不同物种种数的线性增加，会由于对每个物种的竞争建模复杂度上升而导致运行时间的非线性增加。
\item 我们的模型只考虑水分稀缺程度变化的影响及其对植物竞争相应的作用，并假定土壤养分、日照暴露等其他因素对所有植物群落都是恒定的。
\end{itemize}

\section{结论}

我们的植物生长模型提供了证据，支持了最初的观察：生物多样性的提高能够增强植物群落在干旱周期中的存续能力。通过纳入可模拟多轮干旱周期的自定义生成降雨数据，我们回应了问题重述中提出的每一点：

\begin{enumerate}
\item 尽管干旱时期存在种间竞争，多个植物物种仍彼此互利，因为植物总生物量的增加使得因蒸发或排水而造成的水分损失更少，水分保存得更多。
\item 只需两个物种即可从生物多样性现象中获益。这也正是边际收益最大的地方，此后每增加一个物种都出现收益递减。随着植物数目的增加，我们在空间密度图中看到生物量对空间更大的占据。这模拟了不同物种环境生态位的差异如何改善植物覆盖度，从而减少因排水造成的水分损失。
\item 干旱的频率对植物群落总生物量的损害不如其严重程度那么大。植物一般能从高频但低严重度的干旱中恢复，而在严重干旱期间，物种仍存在灭绝或生物量近乎消失的风险。
\item 每个物种的恢复力与抗性特征分别影响该物种的生长率与死亡率。如预期那样，恢复力下降与污染加重都会导致长期生物量减少。
\item 鉴于模型的空间性质，模拟栖息地缩减只是减小可供植物生长的总面积。我们发现，在物种数目相同的情况下竞争加剧，而在高度受限的区域中，物种甚至可以主导整个系统，或完全灭绝。
\end{enumerate}

因此，为确保植物群落的长期存续能力，必须保护并维持生物多样性。植物群落中生物多样性的存在以种间竞争为代价，这因而自然降低了环境中每个单独物种的生物量，但它通过提高净总植物生物量、提高在干旱周期与其他意外胁迫下存活并适应的机会而有所助益。

\section{未来探索}

我们当前的模型假定：无论环境如何，水分都以与当前水量成比例的速率逸出，且对所有生物群落采用相同的比例常数。基于不同生物群落的气候与土壤数据建立一个体系，将使我们能够在不修改核心模型本身的前提下，把模型应用于更多样的条件集合。此外，把模型扩展到刻画山地丘陵等复杂的非平坦地形，可以为判断哪些区域在干旱条件下对植物最具存活性提供更多洞察。

\section*{参考文献}

\begin{enumerate}
\item[{[1]}] A. Kaisermann, F. T. de Vries, R. I. Griffiths, and R. D. Bardgett, ``Legacy effects of drought on plant--soil feedbacks and plant--plant interactions,'' \emph{New Phytologist}, vol. 215, no. 4, pp. 1413--1424, 2017.
\item[{[2]}] S. Direct. Plant communities. [Online]. Available: https://www.sciencedirect.com/topics/agricultural-and-biological-sciences/plant-communities
\item[{[3]}] F. T. de Vries, M. E. Liiri, L. Bj\o rnlund, H. M. Set\"al\"a, S. Christensen, and R. D. Bardgett, ``Legacy effects of drought on plant growth and the soil food web,'' \emph{Oecologia}, vol. 170, pp. 821--833, 2012.
\item[{[4]}] T. S. Writer. Ecological resilience and resistance. [Online]. Available: https://www.thesciencewriter.org/resilience-stories/ecological-resilience-and-resistance
\item[{[5]}] S. Yachi and M. Loreau, ``Biodiversity and ecosystem productivity in a fluctuating environment: The insurance hypothesis,'' \emph{Proceedings of the National Academy of Sciences}, vol. 96, no. 4, pp. 1463--1468, 1999.
\item[{[6]}] A. J. Lotka, ``Contribution to the theory of periodic reactions,'' Jan. 1909.
\item[{[7]}] J. Monod, ``The growth of bacterial cultures,'' \emph{Annual Review of Microbiology}, vol. 3, no. 1, pp. 371--394, 1949. [Online]. Available: https://doi.org/10.1146/annurev.mi.03.100149.002103
\item[{[8]}] H. H. Elsalahy, S. D. Bellingrath-Kimura, C.-L. Ro\ss, T. Kautz, and T. F. D\"oring, ``Crop resilience to drought with and without response diversity,'' \emph{Frontiers in Plant Science}, vol. 11, 2020. [Online]. Available: https://www.frontiersin.org/articles/10.3389/fpls.2020.00721
\item[{[9]}] D. U. Hooper and P. M. Vitousek, ``The effects of plant composition and diversity on ecosystem processes,'' \emph{Science}, vol. 277, no. 5330, pp. 1302--1305, 1997.
\item[{[10]}] T. D., D. Wedin, and K. J., ``Productivity and sustainability influenced by biodiversity in grassland ecosystems,'' \emph{Nature}, vol. 379, pp. 718--720, 1996.
\item[{[11]}] S. Naeem, L. Thompson, and L. S., ``Declining biodiversity can alter the performance of ecosystems,'' \emph{Nature}, vol. 368, pp. 734--747, 1994.
\item[{[12]}] J. A. Sherratt, \emph{Vegetation Stripes in Semi-Arid Environments}. University of Sheffield, 2006.
\item[{[13]}] G. Tumble. What is the effect of the pollution on plants. [Online]. Available: https://greentumble.com/effect-of-pollution-on-plants/
\end{enumerate}

\end{document}
